\begin{afterword}
\summaryhead

The post opens with the Blues and Greens of the Roman Empire, factions that fought and murdered each other for generations. Quoting Procopius and Gibbon on their hatred and their power, it reveals that they were supporters of chariot-racing teams.

It then imagines a future society living in sealed caverns, where two factions formed over whether the word ``cerulean,'' used in an old description of the sky, means blue or green. The factions have since taken opposite positions on taxes, marriage and cosmology. After long violence there is a truce, and some tolerance.

An earthquake opens a passage to the surface, and one member of a party of six climbs up and sees a blue sky. The story branches six ways, depending on who climbed. Aditya, a Blue, resolves to end the truce. Barron, a Green, despairs. Charles, a Blue professor, plans to block the passage. Daria, a Green, painfully accepts that she was wrong. Eddin, a Green, laughs at the pointlessness of it all. Ferris looks around in delight and goes exploring.

\responsehead

The fable is well made, and its six portraits are easy to recognize. What it shows is narrower than its title.

Its one historical claim is stated more firmly than the history allows. The Blues and Greens began as racing clubs, and Procopius, in the passage the post quotes, says the partisans fought for the names and the seats at the games. But historians disagree about how far the factions had become political and religious parties by the sixth century, and Gibbon's next sentence, which the post leaves out, ties the Greens to the sect of Anastasius and the Blues to orthodoxy. ``They were sports fans'' is one side of that debate, offered as the answer.

The test the fable sets its characters is a fact that anyone can see by walking up a corridor. The factions' packages also hold questions of value: which tax to levy, how easy divorce should be. The sky settles none of these. The fable shows how people react when one question in their side's package is settled beyond doubt. It does not show how to reason about the questions a sky cannot settle, such as the tax and the marriage laws, which are the political part of each package.

The fable states no moral, and a reader must infer one from the portraits. The book later supplies it: in ``The Meditation on Curiosity,'' Ferris ``embodies the power of innocent curiosity.'' But the fable gives that response only to Ferris, the one member of the party with no faction and no pronoun. The only partisan who reaches the truth, Daria, does so through grief. So the fable describes the model response and the partisans' failures, but not how a partisan might reach it. Charles's case is also left open. His fear that the discovery will be ``misinterpreted'' is exactly what happens in Aditya's branch, and the fable does not say whether hiding the sky is wrong.

In short: a well-told set of reactions to a fact anyone can check, opened by a historical claim that takes one side of a debate, and silent on the questions of value that the sky cannot settle.
\end{afterword}
